%% methodology.tex
%%

%% ==============================
\section{Methodology}
\label{sec:Methodology}
%% ==============================

\subsection{Synthetic Multi-Site QoS Telemetry}
Real multi-site OneM2M telemetry was not available, so a synthetic generator (\texttt{src/data/qos\_simulator.py}) produces QoS records (RTT, CPU, RAM, success rate, latency, throughput) for 6 simulated IoT edge sites. Each site's traffic is a mixture of the baseline paper's own three scenarios -- uniform, burst, and real-time -- but with a different, skewed mixture per site (80\% dominated by one scenario, 10\% each of the other two), which is what makes this a genuine non-IID federated learning setting: no single site sees a representative sample of the overall traffic distribution. 1,500 records per site are generated per random seed.

\subsection{Training}
For each of 5 seeds (42-46), each site's data is split 80/20 (stratified by label) into a local train/test set. The centralized model pools all sites' training splits and SMOTE-balances the pooled result before fitting. The federated model fits one Random Forest per site on that site's training split only, SMOTE-balanced locally. Both use identical Random Forest hyperparameters (200 trees, max depth 10) so any performance difference reflects the data-sharing strategy, not model capacity. All three variants (centralized, federated, single-site-only) are evaluated on the same pooled, held-out test set (20\% of every site's data, never used in training by any model).

\subsection{Metrics}
\begin{itemize}
  \item \textbf{Accuracy / Macro-F1}: over the full 3-class pooled test set.
  \item \textbf{Critical Recall (class 2)}: recall specifically on the ``Full Offload'' class -- the minority class the baseline paper needed SMOTE to handle (9.64\% of their real data), and the one most consequential to miss in production (a missed Critical event means an overloaded system stays local).
\end{itemize}

All code and results are in \texttt{iot-reliability/}.
